\section{Teilversuch 1}
% Beispielinhalt: alles unterhalb der \section-Zeile kann gelöscht werden.

\blindtext

\begin{figure}[htbp]
  \centering
  \begin{subfigure}{0.45\textwidth}
    \includegraphics[width=\linewidth]{example-image-b}
    \caption{Buchstabe B}
    \label{fig:b}
  \end{subfigure}\hfill
  \begin{subfigure}{0.45\textwidth}
    \includegraphics[width=\linewidth]{example-image-c}
    \caption{Buchstabe C}
    \label{fig:c}
  \end{subfigure}
  \caption{Zwei Teilabbildungen nebeneinander. Lange Beschriftungen werden
    im Blocksatz gesetzt, kurze zentriert}
  \label{fig:bc}
\end{figure}

\Cref{fig:bc} zeigt zwei Buchstaben; \cref{fig:c} lässt sich auch einzeln
referenzieren. Setzt man \cref{eq:dyson} beliebig oft in sich selbst ein,
erhält man angeblich die Werte in \cref{tab:werte}.

\begin{table}[htbp]
  \caption{Zahlen ohne Bedeutung, an der Dezimalstelle ausgerichtet}
  \label{tab:werte}
  \centering
  \begin{tabular}{
      S[table-format=1.2]
      S[table-format=2.2]
      S[table-format=2.1(2)]}
    \toprule
    {$x$ / \unit{\metre}} & {$y$ / \unit{\second}}
      & {$z$ / \unit{\kilo\gram}} \\
    \midrule
    0.25 & 14.47 & 14.1(8)  \\
    0.50 & 28.94 & 29.5(8)  \\
    0.75 & 43.41 & 42.7(9)  \\
    1.00 & 57.88 & 58.6(10) \\
    \bottomrule
  \end{tabular}
\end{table}

\begin{figure}[htbp]
  \centering
  \includegraphics[width=0.7\textwidth]{example-image}
  \caption{Eine breite Platzhalterabbildung}
  \label{fig:d}
\end{figure}

Aus \cref{fig:d} folgt mit etwas Fantasie
$\zeta = \qty{58.1(9)}{\micro\electronvolt\per\tesla}$.

\blindtext
